\chapter{Content Guide}\label{record-040}

\noindent{\footnotesize\color{BookMuted}\path{docs/strategy/content-guide.md}\par}\smallskip

\emph{Voice, tone, and copy doctrine for Bijou interfaces}

\section{Why this exists}\label{record-040-why-this-exists}

In Bijou, words are part of the product.

This is especially true in terminal interfaces, where:

\begin{itemize}
\tightlist
\item
  labels are compact
\item
  hints compete for space
\item
  accessible and pipe modes rely heavily on text
\item
  empty states and confirmations must do real orientation work
\end{itemize}

This guide exists to keep Bijou copy:

\begin{itemize}
\tightlist
\item
  calm
\item
  clear
\item
  localizable
\item
  accessible
\item
  operationally honest
\end{itemize}

\section{Voice}\label{record-040-voice}

Bijou copy should sound:

\begin{itemize}
\tightlist
\item
  direct
\item
  calm
\item
  precise
\item
  unpretentious
\end{itemize}

This applies to both human-facing and agent-facing surfaces.

Avoid copy that feels:

\begin{itemize}
\tightlist
\item
  hyped
\item
  mysterious
\item
  overly cute
\item
  apologetic
\item
  passive when action is needed
\end{itemize}

\section{Tone by surface}\label{record-040-tone-by-surface}

\subsection{Shell chrome}\label{record-040-shell-chrome}

Shell copy should be terse and operational.

Examples:

\begin{itemize}
\tightlist
\item
  \texttt{?\allowbreak{} \allowbreak{}H\allowbreak{}e\allowbreak{}l\allowbreak{}p}
\item
  \texttt{/\allowbreak{} \allowbreak{}S\allowbreak{}e\allowbreak{}a\allowbreak{}r\allowbreak{}c\allowbreak{}h}
\item
  \texttt{F\allowbreak{}2\allowbreak{} \allowbreak{}S\allowbreak{}e\allowbreak{}t\allowbreak{}t\allowbreak{}i\allowbreak{}n\allowbreak{}g\allowbreak{}s}
\end{itemize}

Chrome should not carry essay-length explanation.

\subsection{Empty states}\label{record-040-empty-states}

Empty states should orient.

They should answer:

\begin{itemize}
\tightlist
\item
  what this place is
\item
  why it is empty
\item
  what to do next
\end{itemize}

They should not default to:

\begin{itemize}
\tightlist
\item
  ``Nothing here''
\item
  ``No data''
\item
  ``Oops''
\end{itemize}

\subsection{Confirmations}\label{record-040-confirmations}

Confirmations should state consequence, not merely ask a generic yes/no question.

Good:

\begin{itemize}
\tightlist
\item
  \texttt{Q\allowbreak{}u\allowbreak{}i\allowbreak{}t\allowbreak{} \allowbreak{}B\allowbreak{}i\allowbreak{}j\allowbreak{}o\allowbreak{}u\allowbreak{} \allowbreak{}D\allowbreak{}o\allowbreak{}c\allowbreak{}s\allowbreak{}?}
\item
  \texttt{D\allowbreak{}i\allowbreak{}s\allowbreak{}c\allowbreak{}a\allowbreak{}r\allowbreak{}d\allowbreak{} \allowbreak{}u\allowbreak{}n\allowbreak{}s\allowbreak{}a\allowbreak{}v\allowbreak{}e\allowbreak{}d\allowbreak{} \allowbreak{}f\allowbreak{}i\allowbreak{}l\allowbreak{}t\allowbreak{}e\allowbreak{}r\allowbreak{} \allowbreak{}c\allowbreak{}h\allowbreak{}a\allowbreak{}n\allowbreak{}g\allowbreak{}e\allowbreak{}s\allowbreak{}?}
\end{itemize}

Better body copy:

\begin{itemize}
\tightlist
\item
  what will happen
\item
  what will not happen
\item
  what the safe choice is when needed
\end{itemize}

\subsection{Errors}\label{record-040-errors}

Errors should try to provide:

\begin{itemize}
\tightlist
\item
  what happened
\item
  basis if known
\item
  next action if possible
\end{itemize}

Avoid vague panic language.

\subsection{Notifications}\label{record-040-notifications}

Notifications should be short, factual, and scannable.

They should confirm:

\begin{itemize}
\tightlist
\item
  what changed
\item
  what succeeded
\item
  what needs attention
\end{itemize}

without pretending to be a full explanation surface when they are not.

\subsection{Explainability surfaces}\label{record-040-explainability-surfaces}

Explainability copy should identify:

\begin{itemize}
\tightlist
\item
  source
\item
  rationale
\item
  evidence
\item
  confidence posture
\item
  next action
\end{itemize}

\section{Writing rules}\label{record-040-writing-rules}

\subsection{1. Prefer explicit over clever}\label{record-040-prefer-explicit-over-clever}

Especially in:

\begin{itemize}
\tightlist
\item
  accessible mode
\item
  onboarding copy
\item
  help
\item
  confirmations
\end{itemize}

\subsection{2. Use verb-first action labels}\label{record-040-use-verb-first-action-labels}

Examples:

\begin{itemize}
\tightlist
\item
  \texttt{O\allowbreak{}p\allowbreak{}e\allowbreak{}n\allowbreak{} \allowbreak{}s\allowbreak{}e\allowbreak{}t\allowbreak{}t\allowbreak{}i\allowbreak{}n\allowbreak{}g\allowbreak{}s}
\item
  \texttt{S\allowbreak{}e\allowbreak{}a\allowbreak{}r\allowbreak{}c\allowbreak{}h\allowbreak{} \allowbreak{}c\allowbreak{}o\allowbreak{}m\allowbreak{}p\allowbreak{}o\allowbreak{}n\allowbreak{}e\allowbreak{}n\allowbreak{}t\allowbreak{}s}
\item
  \texttt{D\allowbreak{}i\allowbreak{}s\allowbreak{}m\allowbreak{}i\allowbreak{}s\allowbreak{}s\allowbreak{} \allowbreak{}n\allowbreak{}o\allowbreak{}t\allowbreak{}i\allowbreak{}f\allowbreak{}i\allowbreak{}c\allowbreak{}a\allowbreak{}t\allowbreak{}i\allowbreak{}o\allowbreak{}n}
\end{itemize}

\subsection{3. Avoid unexplained jargon in first-use paths}\label{record-040-avoid-unexplained-jargon-in-first-use-paths}

Terms like:

\begin{itemize}
\tightlist
\item
  lane
\item
  pane
\item
  profile
\item
  lowering
\end{itemize}

should be introduced with care when the user is new.

\subsection{4. Treat hint bars as scarce space}\label{record-040-treat-hint-bars-as-scarce-space}

Hints should be:

\begin{itemize}
\tightlist
\item
  short
\item
  truthful
\item
  current
\end{itemize}

They should not attempt to explain the whole app.

\subsection{5. Write for localization}\label{record-040-write-for-localization}

Avoid:

\begin{itemize}
\tightlist
\item
  English-specific jokes
\item
  brittle idioms
\item
  copy that only works at one exact length
\end{itemize}

\subsection{6. Write for accessible linear reading}\label{record-040-write-for-accessible-linear-reading}

When a surface lowers to accessible mode, the copy should still make sense when read top-to-bottom with no spatial context.

\subsection{7. Write for agents as first-class users}\label{record-040-write-for-agents-as-first-class-users}

Agent-facing copy, labels, and explanations should be:

\begin{itemize}
\tightlist
\item
  explicit
\item
  inspectable
\item
  deterministic
\item
  low on hidden implication
\end{itemize}

\section{AI-specific copy posture}\label{record-040-ai-specific-copy-posture}

When AI is present:

\begin{itemize}
\tightlist
\item
  mark it
\item
  explain why it is present
\item
  explain what the user should do next
\item
  avoid anthropomorphic bluffing
\end{itemize}

See \href{https://github.com/flyingrobots/bijou/blob/cbd70612140e7a450d764b3c7be1c6af878b99bc/docs/strategy/ai-explainability-standard.md}{AI Explainability Standard}.

\section{Relationship to other doctrine}\label{record-040-relationship-to-other-doctrine}

This note should be read alongside:

\begin{itemize}
\tightlist
\item
  \href{https://github.com/flyingrobots/bijou/blob/cbd70612140e7a450d764b3c7be1c6af878b99bc/docs/strategy/bijou-ux-doctrine.md}{Bijou UX Doctrine}
\item
  \href{https://github.com/flyingrobots/bijou/blob/cbd70612140e7a450d764b3c7be1c6af878b99bc/docs/strategy/accessibility-and-assistive-modes.md}{Accessibility and Assistive Modes}
\item
  \href{https://github.com/flyingrobots/bijou/blob/cbd70612140e7a450d764b3c7be1c6af878b99bc/docs/strategy/localization-and-bidirectionality.md}{Localization and Bidirectionality}
\item
  \href{https://github.com/flyingrobots/bijou/blob/cbd70612140e7a450d764b3c7be1c6af878b99bc/docs/strategy/ai-explainability-standard.md}{AI Explainability Standard}
\end{itemize}

\section{Acceptance bar}\label{record-040-acceptance-bar}

This guide is working when:

\begin{itemize}
\tightlist
\item
  shell copy is consistent and calm
\item
  empty states teach instead of merely reporting absence
\item
  confirmations explain consequence
\item
  accessible and localized surfaces have usable prose
\item
  Bijou's tone feels intentional across apps and docs
\end{itemize}

